%% ============================================================
%%  LỜI CẢM ƠN
%% ============================================================
\chapter*{Lời cảm ơn}
\addcontentsline{toc}{chapter}{Lời cảm ơn}
\markboth{LỜI CẢM ƠN}{}

Em xin bày tỏ lòng biết ơn sâu sắc đến quý thầy cô Khoa Mạng máy tính và Truyền thông, Trường Đại học Công nghệ Thông tin, Đại học Quốc gia Thành phố Hồ Chí Minh, những người đã truyền đạt nền tảng kiến thức chuyên môn, phương pháp tư duy khoa học và tinh thần làm việc nghiêm túc trong suốt quá trình học tập. Những bài giảng, những buổi thảo luận và những góp ý học thuật đã tạo nên cơ sở vững chắc để đề tài này được hình thành và phát triển theo hướng bài bản, có hệ thống và phù hợp với định hướng nghiên cứu của ngành.

Đặc biệt, em xin trân trọng cảm ơn ThS. Nguyễn Hữu Quyền, giảng viên hướng dẫn của báo cáo thực tập. Trong suốt quá trình thực hiện đề tài, từ giai đoạn xác định vấn đề nghiên cứu, xây dựng cơ sở lý thuyết, thiết kế hệ thống cho đến triển khai thực nghiệm và hoàn thiện bản thảo, sự hướng dẫn tận tình và những góp ý chi tiết của thầy đã giúp em nhìn nhận vấn đề một cách chặt chẽ và toàn diện hơn. Những buổi trao đổi chuyên môn không chỉ dừng lại ở việc chỉnh sửa nội dung, mà còn định hình phương pháp tiếp cận khoa học, cách tổ chức lập luận và cách trình bày kết quả một cách mạch lạc. Sự nghiêm túc trong học thuật và tinh thần trách nhiệm của thầy là nguồn động lực quan trọng để báo cáo này được hoàn thiện với chất lượng tốt nhất trong khả năng của em.

Bên cạnh đó, em xin gửi lời cảm ơn đến tập thể cán bộ, giảng viên và nhân viên của Khoa đã tạo điều kiện thuận lợi về cơ sở vật chất, tài liệu tham khảo và môi trường học tập.

Em cũng xin trân trọng cảm ơn Trung tâm An ninh mạng đã tạo điều kiện thực tập, cung cấp môi trường làm việc chuyên nghiệp, cơ sở hạ tầng kỹ thuật và những bài toán thực tế có giá trị, giúp em áp dụng kiến thức học thuật vào thực tiễn một cách hiệu quả.

% Em cũng xin bày tỏ lòng biết ơn đối với gia đình, những người luôn âm thầm ủng hộ và tạo điều kiện tốt nhất để việc học tập và nghiên cứu được diễn ra liên tục và ổn định. Sự động viên, tin tưởng và chia sẻ từ gia đình là chỗ dựa tinh thần quan trọng trong những giai đoạn gặp khó khăn về học thuật cũng như áp lực tiến độ.

Cuối cùng, mặc dù vẫn còn một số hạn chế và nhiều hướng có thể tiếp tục cải thiện, báo cáo đã đạt được các mục tiêu đề ra và cung cấp những kết quả ban đầu khả quan cho bài toán nghiên cứu. Những góp ý và nhận xét từ quý thầy cô sẽ là cơ sở quan trọng để tác giả tiếp tục hoàn thiện và phát triển hướng nghiên cứu này trong tương lai.

\vspace{2cm}
\begin{flushright}
\textit{TP. Hồ Chí Minh, tháng 06 năm 2026}\\[0.5cm]
% \textbf{Tác giả}
\end{flushright}

\pagebreak
